% 学習結果説明（手順5b：最終候補 B_1・C_1 の MCMC。用語集と図の見方）
% results フォルダで「lualatex 学習結果説明_手順5b.tex」を2回実行すると PDF を作り直せる（TeX Live の LuaLaTeX を使用）。
% 図は results/step5b_mcmc/ の PNG を相対パスで読み込む。
\documentclass[paper=a4paper,fontsize=10.5pt,line_length=44zw,number_of_lines=38]{jlreq}
\usepackage[haranoaji]{luatexja-preset}
\usepackage{amsmath}
\usepackage{graphicx}
\usepackage{booktabs}
\usepackage{tabularx}
\usepackage{enumitem}
\usepackage{float}
\usepackage{xcolor}
\usepackage[hidelinks,bookmarksnumbered]{hyperref}
\graphicspath{{step5b_mcmc/}}

\newcommand{\file}[1]{\texttt{\detokenize{#1}}}
\newcommand{\gd}{\dot{\gamma}}
\newcommand{\Tff}{T_{\mathrm{ff}}}
\newcommand{\Tset}{T_{\mathrm{set}}}
\newcommand{\Eh}{E_{h}}
\newcommand{\Ev}{E_{v}}
\newcommand{\PI}{P3\_1}
\newcommand{\PII}{P3\_2}
\newenvironment{termlist}{\begin{itemize}[label=・,leftmargin=1\zw,labelwidth=1\zw,labelsep=0pt,itemsep=0.4ex,topsep=0.6ex,parsep=0pt]}{\end{itemize}}
\newcommand{\term}[1]{\item \textbf{#1}：}
\newenvironment{points}{\begin{itemize}[label=$\triangleright$,leftmargin=1.5\zw,itemsep=0.3ex,topsep=0.5ex,parsep=0pt]}{\end{itemize}}
\newcommand{\pt}[1]{\item \textbf{#1}：}
\renewcommand{\arraystretch}{1.25}

\title{樹脂粘度モデル開発　手順5bの学習結果の説明\\[0.6ex]{\large 最終候補（B\_1・C\_1）の MCMC：用語集とグラフの見方（機械学習・統計を初めて学ぶ人向け）}}
\author{}
\date{@@DATE@@}

\begin{document}
\maketitle
{\small\tableofcontents}

% =====================================================================
\section{この文書について}

この文書は、\file{results/step5b_mcmc/} にある手順5bの結果（図と表）を、機械学習や統計を初めて学ぶ人でも読めるように説明したものです。手順5bは、手順5で決まった最終候補のモデルについて、ハイパーパラメータと温度のずらし量の\textbf{不確かさ}を、MCMC という方法で調べた手順です。

\begin{points}
\pt{対象の結果}\file{results/step5b_mcmc/}：図 @@N_FIG@@ 枚、CSV @@N_CSV@@ 個、事後サンプル @@N_POST@@ 個。図は、すべて \file{results/step5b_mcmc/highres/} にも 600 dpi で保存されます（ほかに、せん断速度 ― 粘度のグラフ @@N_HR@@ 枚）。
\pt{読み方}2章で何を調べたかとモデルの構成をつかみ、分からない言葉は3章の用語集で調べてください。4章で図の見方、5章で表の見方、6章で結論と、これからの選択肢を説明します。\textbf{結論だけを知りたいときは、6章から読んでください}。
\pt{元のノート}\file{01_機械学習コードフォルダ} の \file{手順5b_最終候補のMCMC.ipynb} です。実行は約 @@RUN_MIN@@ 分です。
\pt{記号表}記号・モデル名・列名の一覧は、\file{学習結果説明_記号表.pdf} にあります。
\pt{作り直し方}\file{results} フォルダで \texttt{lualatex 学習結果説明\_手順5b.tex} を2回実行すると、この PDF を作り直せます。
\pt{ノートの判定}ノートの各節の最後に、\file{>> 判定：OK}（想定どおり）か \file{>> 判定：★要確認}（想定と違う、または注意が要る点）が出ます。今回は、OK が @@N_OK@@ 個、★ が @@N_STAR@@ 個でした。★ はエラーではなく、結果を読むときに気をつける点の目印です。6章に、★ の意味を書きました。
\end{points}

% =====================================================================
\section{何を調べたか（概要）}

\subsection{最尤推定と MCMC の違い}
ガウス過程回帰には、データから決める数（ハイパーパラメータ）があります（関数の大きさ、入力ごとの滑らかさ、ノイズの大きさなど）。手順5までは、データに最もよく合う 1 組の値を探す\textbf{最尤推定}で決めていました。この方法は、「その 1 組の値が正しい」とみなして予測します。

実際には、データが限られているので、ほかの値の組でも、ほとんど同じくらいデータに合うことがあります。\textbf{MCMC}（マルコフ連鎖モンテカルロ法）は、「データに合う度合いに応じた確率」で、値の組を何万個も取り出し（サンプリング）、値の\textbf{不確かさ}ごと予測に取り込みます。取り出した値の分布を\textbf{事後分布}と呼びます。事後分布が狭ければ、その値はデータでよく決まっています。広ければ、データだけでは決まっていません。

今回のモデルでは、温度のずらし量（横方向 $\Eh$、縦方向 $\Ev$）も、ハイパーパラメータと同時にサンプリングしました（ユーザー決定、10 月 14 日）。手順5では、ずらし量は、格子の中の最良の値に固定していました。

\subsection{モデルの構成}
手順5と同じモデルです。次の表に、学習データ、入力、平均関数、カーネル、事前分布をまとめます。

\begin{table}[H]
\caption{モデルの構成（B\_1：構造 B で厚み $H$ を残す、C\_1：構造なし。C\_1 は手順3b・3d の R3 と同じ）}\label{tab:model}
\small
\begin{tabularx}{\linewidth}{@{}l>{\raggedright\arraybackslash}X@{}}
\toprule
項目 & 内容 \\
\midrule
学習データ & 外れ値なし。問題データ一覧の「除外」（P01・P03・P08）を除く。1 行が 1 ショット。5data 版：5data（PP 460 行、粉砕材混合 PP 459 行）。Alldata 版：Alldata（PP 1234 行、粉砕材混合 PP 560 行）。S\_T の学習は、それぞれ 190・230 ℃ だけ \\
目的変数 & $\ln\eta$（$\eta$：見かけの粘度 [Pa$\cdot$s]）。学習データの平均 $y_0$ を引く（固定の定数。自動の標準化はしない） \\
入力：C\_1 & 条件平均の $\ln\gd$、$\ln H$、$\Tset$（摂氏のまま）の 3 列。学習データの平均と標準偏差で標準化 \\
入力：B\_1 & $u=\ln\gd$（条件平均）$+\ln a(\Tset;\Eh)$、$w=\ln\mathrm{Fo}=-(\ln\gd+3\ln H)$ の 2 列。出力から $\ln a(\Tset;\Ev)$ を引く。$\ln a(T;E)=(E/R)(1/T-1/T_{\mathrm{ref}})$、$T_{\mathrm{ref}}=483.15$ K、$R=8.314$ \\
使わない入力 & \PI・\PII・$\Tff$、ショットごとの $\gd$ \\
平均関数 & 明示的なものはない（ガウス過程の平均はゼロ）。実質的な平均の役目をするのは、定数 $y_0$、B\_1 の縦シフト $\ln a(\Tset;\Ev)$、カーネルの線形項（学習範囲の外へ一次式で延ばす役目） \\
カーネル & $C_1\cdot\mathrm{Matern}(\nu=5/2,\ \text{入力ごとの長さスケール})+C_2\cdot\mathrm{DotProduct}(\sigma_0)+$ノイズ（分散 $\sigma_n^2$） \\
事前分布：カーネル & 手順3c と同じ。探索範囲の中で、対数の尺度で一様（$C_1$・$C_2$：$10^{-3}$〜$10^{3}$、長さスケール：$10^{-2}$〜$10^{2}$、$\sigma_0$：$10^{-5}$〜$10^{3}$、$\sigma_n^2$：$10^{-6}$〜$10$）。範囲の外は確率 0 \\
事前分布：ずらし量 & 範囲の中で一様（$\Eh$：$-40$〜$40$、$\Ev$：$-20$〜$40$ kJ/mol。手順5のプロファイル尤度の格子と同じ範囲） \\
関数の事前分布 & 関数 $f\sim$ ガウス過程（平均 0、共分散は上のカーネル）。ノイズ $\varepsilon\sim N(0,\sigma_n^2)$ \\
\bottomrule
\end{tabularx}
\end{table}

\subsection{計算の工夫（条件ごとにまとめた尤度）}
入力は条件平均なので、同じ条件のショットは、入力が同じです。そこで、尤度（データがそのモデルでどれだけありそうか）を、条件ごとの平均・ショット数・条件内の平方和を使って、\textbf{厳密に}書き直しました（ショットごとに計算した結果と同じ値になります）。条件の数（93〜129）の大きさの行列だけで済むので、1 回の計算が約 1 ミリ秒です。Alldata 版（PP は 1234 行）でも同じ速さです。ノートの最初で、scikit-learn のショットごとの GPR と、対数周辺尤度・予測が一致することを確かめています（差は $10^{-9}$ 程度）。

\subsection{手順の流れ}
\begin{enumerate}[itemsep=0.3ex,topsep=0.5ex]
\item 最尤推定を 2 通り作る。「MLE（E 固定）」は、ずらし量を手順5の最良値に固定して、カーネルだけを推定（手順5と同じ設定）。「MLE（同時）」は、ずらし量も一緒に推定。
\item 4 本のチェーンで、MCMC のサンプリングを行う（調整 5000 回は捨て、本番 20000 回。収束しなければ回数を増やす）。
\item 事後サンプル 200 個の予測分布を等しい重みで混ぜた「混合分布」で、検証データを予測し、最尤推定と比べる。
\end{enumerate}
対象は、材料 2 $\times$ モデル 2 $\times$ 学習データ 4（5data 版：S\_V・S\_T、Alldata 版：S\_T・ALL）= 16 組です。

% =====================================================================
\section{用語集}
\begin{termlist}
\term{MCMC}2.1 節。値の組を、データに合う度合いに応じた確率で、何万個も取り出す方法。ここでは、自作の「適応型ランダムウォーク・メトロポリス法」（手順2・3c と同じ）。
\term{事後分布・事前分布}事前分布：データを見る前の、値についての考え（ここでは、決めた範囲の中で一様）。事後分布：データを見たあとの、値の分布。
\term{95\% 区間}事後サンプルを小さい順に並べたときの、2.5\% から 97.5\% の範囲。「値がこの範囲にある確からしさが 95\%」と読む。
\term{チェーン}別々の出発点から歩かせた、一続きのサンプル（4 本）。
\term{調整期間・本番}調整期間（5000 回）：歩幅などを合わせる期間で、捨てる。本番（20000 回）：結果に使う。
\term{採択率}提案した次の点に、実際に移れた割合。目標 0.234。
\term{R-hat}チェーンどうしがそろっているかの指標。1.01 以下が目安。
\term{ESS（有効サンプル数）}隣り合うサンプルは似ているので、独立なサンプル何個分にあたるか。bulk は分布の中心、tail は裾。400 以上が目安。
\term{トレース}サンプルの推移を、反復の順に描いた図。毛虫のように、一定の範囲を行ったり来たりしていれば、よく混ざっている。
\term{MLE（E 固定）・MLE（同時）}最尤推定。E 固定：ずらし量を手順5の最良値に固定。同時：ずらし量も推定。
\term{混合分布}事後サンプルごとの予測分布を、同じ重みで混ぜた分布。ハイパーパラメータとずらし量の不確かさが、予測の幅に入る。
\term{範囲の端への張り付き}事後分布の 95\% 区間が、事前分布の範囲の端に届くこと。そのパラメータは、データだけでは決まらない（範囲の外の値でも、データに合うかもしれない）。
\term{長さスケール}入力がどれだけ変わると出力が大きく変わるか。大きいほど、その入力に対して滑らか。上限（100）に届くと、その入力に対しては、ほぼ直線（線形項が担当）。
\term{線形項の $\sigma_0$}線形項の定数部分の大きさ。予測にほとんど効かないので、データでは決まりにくい。
\term{RMSE・被覆率・NLPD}手順3a・3b の用語集と同じ（予測の平均的な誤差、95\% 予測区間に実測が入った割合、当たり具合と区間の適切さをまとめた指標）。
\term{予測区間と信頼区間}予測区間：新しく測った値が入る範囲（ノイズ込み）。ずらし量の事後区間や、設定温度の効果の区間は、平均の不確かさだけ。
\end{termlist}

% =====================================================================
\section{図の見方}

@@FIGURES@@

% =====================================================================
\section{表の見方}

@@TABLES@@

% =====================================================================
\section{結論と、これからの選択肢}

@@CONCLUSION@@

\end{document}
